\section{Scalar normalisation}
\label{sec:finite-part}

Throughout this section, \(A\) denotes a primitive non-negative real
matrix with Perron root \(\lambda>1\). The symbols \(p_n(A)\) denote the
algebraic M\"obius coefficients of \eqref{eq:primitive-mobius}; they are
primitive orbit counts only in the integer multigraph case fixed in
\Cref{sec:preliminaries}. For the finite-group results, the scalar role is the
trivial-channel finite part used in
\Cref{thm:renormalized-class-mertens,thm:renormalized-frobenius-product} and,
after \Cref{cor:strict-gap-specialization}, in the strict-gap constants.
This section records the finite-part formula and the scalar projection fact
needed for that normalisation.  Fixed-matrix cyclic-lift and root-of-unity
records appear only as comparison records for this scalar layer.

\subsection{Finite parts and reduced determinants}

This subsection is classical scalar infrastructure rather than a new
contribution.  It collects the Perron-pole normalisation of the SFT zeta
function from \Cref{prop:prelim-trace-primitive} and the corresponding
dynamical Mertens finite part in the form needed by the finite-group
sections; compare \cite{BowenLanford1970Zeta,Parry1983PrimeOrbit,Sharp1991Mertens,LindMarcus1995,ParryPollicott1990Zeta}.  We keep the
short proofs to make the constants \(C(A)\), \(\FP(A)\), and
\(\mathcal M(A)\) unambiguous.

\begin{lemma}[Classical scalar normalisation at the Perron pole]
  \label{lem:finite-part-primitive-closed-form}
  Let \(A\) be a primitive non-negative matrix with Perron root
  \(\lambda>1\). This is a scalar \S\ref{sec:finite-part} statement: the
  quantities \(p_n(A)\) are the algebraic M\"obius coefficients and become
  literal primitive orbit counts only in the integer multigraph case. Then
  \begin{equation}\label{eq:finite-part-determinant-mobius}
    \FP(A)
    =
    \log C(A)
    +\sum_{m\ge 2}\frac{\mu(m)}{m}
      \log\det(I-\lambda^{-m}A)^{-1}.
  \end{equation}
  Equivalently,
  \begin{equation}\label{eq:finite-part-primitive-closed}
    \FP(A)
    =
    \log C(A)
    +\sum_{n\ge 1}p_n(A)
      \bigl(\lambda^{-n}+\log(1-\lambda^{-n})\bigr).
  \end{equation}
  Both series converge absolutely. In particular, if
  \[
    \psi(x):=-x-\log(1-x)=\sum_{k\ge 2}\frac{x^k}{k},
    \qquad 0<x<1,
  \]
  then
  \begin{equation}\label{eq:finite-part-gap-series}
    \log C(A)-\FP(A)
    =
    \sum_{n\ge 1}p_n(A)\psi(\lambda^{-n}).
  \end{equation}
  Thus the scalar Mertens normalisation is
  \(\mathcal M(A)=\gamma+\FP(A)\).  The product normalisation
  \(\gamma+\log C(A)\) is obtained only after adding the scalar product tail
  \(\sum_{r\ge2}\Pi_{\mathbf1}(\lambda^{-r})/r\), as recorded in
  \Cref{lem:scalar-normalisation-two-stage}.
\end{lemma}

\begin{proof}
  This is the standard scalar Euler-product calculation, recorded here
  only to fix the normalisation. Put
  \(\zeta_A(z)=\det(I-zA)^{-1}\). For \(|z|<\lambda^{-1}\),
  \[
    \log\zeta_A(z)
    =
    \sum_{n\ge 1}\frac{\Tr(A^n)}{n}z^n
    =
    \sum_{\ell\ge 1}p_\ell(A)\sum_{k\ge 1}\frac{z^{k\ell}}{k}.
  \]
  M\"obius inversion in the repetition variable gives
  \[
    \sum_{m\ge 1}\frac{\mu(m)}{m}\log\zeta_A(z^m)
    =
    \sum_{\ell\ge 1}p_\ell(A)
      \sum_{q\ge 1}
        \frac{z^{q\ell}}{q}\sum_{m\mid q}\mu(m)
    =
    \sum_{\ell\ge 1}p_\ell(A)z^\ell.
  \]
  With \(z=r/\lambda\), this says
  \[
    P_A(r)
    =
    \log\zeta_A(r/\lambda)
    +\sum_{m\ge 2}\frac{\mu(m)}{m}
      \log\zeta_A\bigl((r/\lambda)^m\bigr).
  \]
  Since
  \[
    \log\zeta_A(r/\lambda)+\log(1-r)
    =
    \log\bigl((1-r)\zeta_A(r/\lambda)\bigr)
    \longrightarrow \log C(A)
  \]
  as \(r\uparrow 1\), it remains to pass to the limit in the \(m\ge2\)
  sum. This is justified by
  \[
    \log\zeta_A(\lambda^{-m})
    =
    \sum_{j\ge 1}\frac{\Tr(A^j)}{j}\lambda^{-mj}
    =
    O(\lambda^{1-m}),
    \qquad m\ge 2,
  \]
  so the determinant series is absolutely summable. Letting
  \(r\uparrow1\) gives \eqref{eq:finite-part-determinant-mobius}.

  For the primitive-orbit form, write the same Euler product at
  \(z=r/\lambda\) as
  \[
    \log\zeta_A(r/\lambda)
    =
    P_A(r)
    +\sum_{n\ge 1}p_n(A)\sum_{k\ge 2}
      \frac{r^{kn}\lambda^{-kn}}{k}.
  \]
  By \Cref{lem:primitive-orbit-asymptotic},
  \(\abs{p_n(A)}=O(\lambda^n/n)\), so the \(k\ge2\) tail is absolutely
  summable at \(r=1\). Hence
  \[
    \FP(A)
    =
    \log C(A)
    -\sum_{n\ge 1}p_n(A)\sum_{k\ge 2}\frac{\lambda^{-kn}}{k}.
  \]
  Finally,
  \[
    -\sum_{k\ge 2}\frac{x^k}{k}=x+\log(1-x),
    \qquad |x|<1,
  \]
  with \(x=\lambda^{-n}\), which proves
  \eqref{eq:finite-part-primitive-closed}. The displayed gap formula is
  the same identity rewritten with \(\psi\).
\end{proof}

\begin{proposition}[Classical reduced determinant identity]
  \label{prop:finite-part-residue-reduced-determinant}
  This is the standard Perron-pole reduced-determinant calculation for a
  finite SFT zeta function; see
  \cite{BowenLanford1970Zeta,LindMarcus1995,ParryPollicott1990Zeta}.  It
  is included only to fix the notation \(C(A)\) used in the scalar
  normalisation.
  Let \(\mu_2,\dots,\mu_d\) be the non-Perron eigenvalues of \(A\). Then
  \[
    C(A)
    =\prod_{j=2}^d\left(1-\frac{\mu_j}{\lambda}\right)^{-1}
    =\det\nolimits'(I-A/\lambda)^{-1}.
  \]
\end{proposition}

\begin{proof}
  Factor
  \[
    \det(I-zA)
    =(1-\lambda z)\prod_{j=2}^d(1-\mu_j z).
  \]
  Multiplying by \(1-\lambda z\) and letting \(z\to\lambda^{-1}\) gives
  the claim.
\end{proof}

\begin{corollary}[Classical scalar orbit-Mertens asymptotic]
  \label{cor:finite-part-mertens-asymptotic}
  \leavevmode\par\noindent
  This is the scalar symbolic-dynamics Mertens statement in the normalisation
  used below.  In Sharp's closed-orbit normalisation the factor is written
  with the exponential weight determined by the entropy; here the entropy is
  \(\log\lambda\), so the weight is \(\lambda^{-n}\).  Compare the classical
  orbit-Mertens literature
  \cite{Sharp1991Mertens,ParryPollicott1990Zeta}; the proof below derives the
  stated constant from the finite-matrix Perron expansion and the scalar
  finite part, so that the later finite-group sections use an explicit
  one-sided SFT convention.
  Let \(A\) be a primitive non-negative matrix with Perron root
  \(\lambda>1\). Then
  \[
    \sum_{n\le N}\frac{p_n(A)}{\lambda^n}
    =\log N+\mathcal M(A)+o(1)
    \qquad (N\to\infty),
  \]
  where \(\mathcal M(A)=\gamma+\FP(A)\).
\end{corollary}

\begin{proof}
  By \Cref{lem:primitive-orbit-asymptotic}, there is a number
  \(0<\vartheta<1\) such that
  \[
    \frac{p_n(A)}{\lambda^n}
    =
    \frac{1}{n}+O\!\left(\frac{\vartheta^n}{n}\right).
  \]
  Hence
  \(\sum_{n\ge1}\abs{p_n(A)\lambda^{-n}-1/n}<\infty\). For
  \(0\le r<1\), the defining series
  \(P_A(r)=\sum_{n\ge1}p_n(A)\lambda^{-n}r^n\) and
  \(-\log(1-r)=\sum_{n\ge1}r^n/n\) give
  \[
    P_A(r)+\log(1-r)
    =
    \sum_{n\ge1}\left(\frac{p_n(A)}{\lambda^n}-\frac1n\right)r^n .
  \]
  The absolute convergence just proved permits termwise passage to the
  limit \(r\uparrow1\). By the definition of \(\FP(A)\),
  \[
    \sum_{n\ge1}\left(\frac{p_n(A)}{\lambda^n}-\frac1n\right)
    =\FP(A).
  \]
  Therefore
  \[
    \sum_{n\le N}\frac{p_n(A)}{\lambda^n}
    =H_N+
      \sum_{n\le N}\left(\frac{p_n(A)}{\lambda^n}-\frac1n\right)
    =H_N+\FP(A)+o(1),
  \]
  because the tail of the absolutely convergent correction series tends to
  zero. Finally \(H_N=\log N+\gamma+o(1)\), so the constant term is
  \(\gamma+\FP(A)=\mathcal M(A)\).
\end{proof}

\subsection{Scalar projection interface}

The remaining scalar material is bookkeeping for the finite-group determinant
and product results: it supplies projection data for the constant
class-function component. By
\Cref{thm:class-product-determinant-rigidity}, the
non-trivial class-product coordinates are the fixed-label Peter--Weyl Perron
values. The canonical cyclic lift
\(A^{\langle q\rangle}=A\otimes P_q\) from \Cref{eq:cyclic-lift-def} and the
root-of-unity layers of \(F_A\) are fixed-matrix scalar comparison records.
The following support statement isolates their normalisation role for
\Cref{thm:renormalized-class-mertens,thm:renormalized-frobenius-product} and
the strict-gap specialization: they give the constant class-function
component exactly.

\begin{theorem}[Support projection: scalar normalisation and no descent]
  \label{thm:finite-part-scalar-interface-no-descent}
  Fix a primitive non-negative matrix \(A\) with Perron root
  \(\lambda>1\), a finite group \(G\), and a non-empty family
  \(\mathscr T\subset\mathcal T_{\mathrm{gap}}(A,G)\) of cocycles
  satisfying the twisted spectral gap; here
  \(\mathcal T_{\mathrm{gap}}(A,G)\) is the set of cocycles
  \(E_A\to G\) with \(\eta(\tau)<\lambda\). For
  \(\tau\in\mathscr T\), let
  \[
    \mathcal M^\#_{A,\tau;G}(C):=\frac{|G|}{|C|}\mathcal M_C(A,\tau;G)
  \]
  and let \(\mathcal B_{A,\tau;G}\) be the normalised primitive class
  profile of \Cref{def:normalised-primitive-class-profile}. Then
  \begin{equation}\label{eq:finite-part-scalar-interface-splitting}
    \mathcal M^\#_{A,\tau;G}
    =
    \mathcal M(A)\mathbf 1_G+
    \mathcal B_{A,\tau;G},
  \end{equation}
  and the first summand is the entire contribution determined solely by
  the scalar matrix \(A\).

  More precisely, let \(\mathcal R(A)\) be any fixed-matrix scalar record,
  including any canonical cyclic-lift record depending only on \(A\).
  The record \(\mathcal R(A)\)
  descends to a single normalised primitive class profile on
  \(\mathscr T\) if and only if the following equivalent conditions hold:
  \begin{enumerate}[label=\textup{(\roman*)}]
    \item \(\mathcal B_{A,\tau;G}\) is independent of
      \(\tau\in\mathscr T\);
    \item for every non-trivial \(\varrho\in\Irr(G)\), the Perron-point
      primitive channel \(\Pi_\varrho^\tau(\lambda^{-1})\) is independent
      of \(\tau\in\mathscr T\).
  \end{enumerate}
  Consequently, if two gap cocycles over the same matrix \(A\) have
  distinct non-trivial Perron-point primitive channels, then all scalar
  finite-part and fixed-matrix cyclic-lift records agree for the two
  cocycles, but their normalised primitive class profiles differ.
\end{theorem}

\begin{proof}
  By \Cref{prop:class-mertens-deviations}, for every conjugacy class
  \(C\subset G\),
  \[
    \mathcal M_C(A,\tau;G)
    =
    \frac{|C|}{|G|}\mathcal M(A)+\Delta_C(A,\tau;G)
  \]
  and
  \[
    \Delta_C(A,\tau;G)
    =
    \frac{|C|}{|G|}
      \sum_{\varrho\neq\mathbf1}
        \chi_\varrho^\ast(C)\,\Pi_\varrho^\tau(\lambda^{-1}).
  \]
  Multiplying the first identity by \(|G|/|C|\) and using the definition
  of \(\mathcal B_{A,\tau;G}\) gives
  \eqref{eq:finite-part-scalar-interface-splitting}. The displayed
  Peter--Weyl expansion contains no trivial character term, so the scalar
  projection of \(\mathcal M^\#_{A,\tau;G}\) is exactly
  \(\mathcal M(A)\mathbf1_G\), while all non-scalar information lies in
  \(\mathcal B_{A,\tau;G}\).

  Every record \(\mathcal R(A)\) named in the statement is a function of
  the single matrix \(A\); it is therefore constant as \(\tau\) varies in
  \(\mathscr T\). Such a fixed record can assign one normalised primitive
  class profile to the whole family exactly when the actual profiles
  \(\mathcal B_{A,\tau;G}\) are all equal. Finally, the irreducible
  characters form a basis of the class-function space, so equality of the
  expansions
  \[
    \mathcal B_{A,\tau;G}(C)
    =
    \sum_{\varrho\neq\mathbf1}
      \chi_\varrho^\ast(C)\,\Pi_\varrho^\tau(\lambda^{-1})
  \]
  for all classes \(C\) is equivalent to equality of every coefficient
  \(\Pi_\varrho^\tau(\lambda^{-1})\) with
  \(\varrho\neq\mathbf1\). Applying this criterion to a two-point family
  gives the final obstruction statement.
\end{proof}


\begin{remark}[Fixed-matrix scalar comparison data and the trivial channel]
  \label{rem:fixed-matrix-scalar-interface-only}
  This is a support statement for the Perron-boundary finite-group theorem
  package and its strict-gap specialization.  In the setting of
  \Cref{thm:finite-part-scalar-interface-no-descent},
  \eqref{eq:finite-part-scalar-interface-splitting} writes the normalised
  profile as the sum of the constant class function
  \(\mathcal M(A)\mathbf1_G\) and the non-trivial Peter--Weyl part
  \(\mathcal B_{A,\tau;G}\).  Thus scalar finite-part formulas and the fixed
  cyclic-lift scalar quantities enter the finite-group theorem chain as
  comparison data for the trivial-channel normalisation
  \(\mathcal M(A)\).  The non-trivial class profile on a family
  \(\mathscr T\) is determined exactly in the independence case characterized in
  \Cref{thm:finite-part-scalar-interface-no-descent}.
\end{remark}

\noindent
The cyclic-lift quantities used here have only this comparison role: scalar
data identify a fixed scalar model and are constant as the finite-group
cocycle varies over a fixed base matrix, while the non-trivial
branch-compatible product coordinates are the fixed-label Peter--Weyl values
studied in the finite-group sections.  The product-side finite-state
consequence is recorded in \Cref{cor:class-product-hyperplane-stability}.

